% Nouveau EGA V (cote 25): archive cote 25, site page 121 (PDF page 122).
% typescript. Draft compared: none.
% AI-transcribed by Claude Opus 5.5 (claude-opus-5-5, adaptive thinking, effort high), Message Batches API, batch ega25b (msgbatch_01A7n5YuXubUPNK1zsepi8RG), 2026-11-13; not yet checked by eye.
\scanpage{25/121}
\mspage{24}

sera automatiquement \deleted{de Cartier}\inserted{diviseur} relatif pour la base $\overset{\hand{\vee}}{P}$ , comme il
résulte du fait que U contient les éléments associés au $\underline{O}_{X_s}$ , s $\in$ S,
condition qui est stable par tout changement de base et en particulier
par S'=$\overset{\hand{\vee}}{P} \longrightarrow$ S. Or il résulte aussitôt du lemme \inserted{4.2.} plus haut que D' se
prolonge en un diviseur \deleted{de Cartier} D \uline{sss} $\underline{N}'$ se prolonge en un Module
inversible, ou encore $i'_{\hand{*}}(\underline{N}')$ est inversible. \deleted{Donc} \over{i}{I}l faudrait rédiger
mieux la fin de la démonstration en termes de cond néc et suff (sans y
aller en deux fois comme je l'ai fait), et dégager aussi \inserted{d'abord} le
\margin{\unsure{voir aussi} \S 20.\illegible{1} et 20.11.}\uline{Corollaire} \inserted{3.4.} \uline{au lemme}\hand{3.2}. Supposons g:X$\to$S plat loc de présentation finie,
soit Z une partie fermée de X telle que x $\in$ Z implique prof$_x$ $\underline{O}_{X_s} \geq 2$
\inserted[soit U=X-Z et i:U$\to$X l'immersion c an.]
(où s = g(x)), \xout{Soit $\underline{L}$ un Module loc libre sur U = X-Z, alors $\underline{L}'$}
\xout{est isomorphe à la restriction} Pour tout Module loc libre \inserted[de t.f] $\underline{L}$ sur X,
considérons sa restriction $\underline{L}' = \underline{L}|U$ . Alors \inserted[a)] le foncteur $\underline{L} \rightsquigarrow \underline{L}'$ est
pleinement fidèle, et pour \xout{\unsure{quelque} Module loc libre $\underline{L}$} tout $\underline{L}$, \xout{xxxx}
l'hom canonique $\underline{L} \to i_{\hand{*}}(\underline{L}')$ est un isomorphisme. \inserted[Pour que $\underline{L}$ soit de rang n, il f et s que $\underline{L}'$ le soit.] \hand{(} b) Soit $\underline{L}'$ un
Module loc libre sur X, alors $\underline{L}'$ est isomorphe à la restriction d'un
Module loc libre $\underline{L}$ sss $i_{\hand{*}}(\underline{L}')$ est loc libre. c) Supposons que $\underline{L}'$ soit
le Module inversible associé à un diviseur \deleted{de Cartier} D'\hand{/}sur U. Alors
la condition envisagée dans b) est aussi néc et suffisante pour que D'
soit la restriction d'un diviseur \deleted{de Cartier} D sur X, \inserted[lequel sera alors unique,]\inserted{et égal à l'adhérence schématique de D' dans X.} \hand{|} Pour que D' soit
diviseur \deleted{de Cartier} relatif par rapport à S, il f. et s. \inserted{\unsure{que}} D le soit.

On utilise simplement le fait que tout $\underline{L}$ satisfait aux \xout{xxx}
hypothèses énoncées pour $\underline{F}$ dans le lemme\hand{3.2}...

\uline{Corollaire} \inserted{3.5.} Supposons \inserted{les anneaux locaux de X factoriels, (par exemple} X régulier\hand{)} \deleted{(par exemple S régulier, et X lisse sur
S)}, Alors l'application $\emptyset \to$ \unsure{$\emptyset'$|$\emptyset$} est bijective. En particulier, si X est
un préschéma régulier loc de type fini sur un corps k, et P un fibré
projectif sur k, il y a correspondance biunivoque entre l'ensemble $\emptyset$ des

\margin{121}
